% Compile twice with XeLaTeX. No shell escape required.
\documentclass[aspectratio=169,10pt]{beamer}
\usepackage[UTF8,fontset=fandol]{ctex}
\usepackage{amsmath,booktabs}
\usetheme{Madrid}
\usecolortheme{seahorse}
\setbeamertemplate{navigation symbols}{}
\setbeamertemplate{footline}[frame number]
\title{用投影解释最小二乘}
\author{雪落长安无声}
\date{2026年11月7日}
\begin{document}
\begin{frame}\titlepage\end{frame}
\begin{frame}{问题：拟合值为什么是投影？}
给定满列秩设计矩阵 $X\in\mathbb R^{n\times p}$ 与响应 $y$，求
\[
 \min_b\|y-Xb\|_2^2.
\]
\begin{block}{本页结论}
最优残差与 $X$ 的每一列正交。
\end{block}
\end{frame}
\begin{frame}{推导：从正规方程得到投影矩阵}
\begin{align*}
 X^{\mathsf T}(y-X\hat b)&=0,\\
 \hat b&=(X^{\mathsf T}X)^{-1}X^{\mathsf T}y,\\
 \hat y&=Hy,\quad H=X(X^{\mathsf T}X)^{-1}X^{\mathsf T}.
\end{align*}
\pause
验证 $H^{\mathsf T}=H$、$H^2=H$，即可确认正交投影。
\end{frame}
\begin{frame}{核对：结论的条件与下一步}
\begin{itemize}
\item 公式要求 $X$ 满列秩；计算中优先使用 QR 或 SVD。
\item 投影关系是代数结论；统计推断另需误差模型。
\item 将完整证明放入论文，答辩每页突出一项结论。
\end{itemize}
\end{frame}
\end{document}
